\documentclass[12pt, a4paper]{academics}

\academicsCoverLabel{解题报告}
\academicsCourse{算法设计与分析}{Design and Analysis of Algorithms}
\academicsReport{线段树与树状数组}{2026 年 6 月 6 日}
\academicsArthur{华博文}{19240212}

\begin{document}

\makeacademicscover

\section{引言}

区间查询与修改是算法竞赛中最为常见的问题类型之一。当数据规模达到 $10^5$ 量级时，朴素的 $\mathrm{O}(nq)$ 扫描策略已完全不可行——我们需要一种能够在对数时间内回答区间查询、支持单点或区间修改的数据结构。线段树与树状数组正是为解决此类问题而生的两类经典工具。

本报告以「区间操作」为主线，选取四道难度递进、思想关联的题目，系统展示从基础树状数组到可持久化线段树的设计范式与优化图景。题一从树状数组的 $\mathrm{lowbit}$ 原理出发，建立前缀和差分的基本模型，实现单点修改与区间查询在 $\mathrm{O}(\log n)$ 时间内完成。题二引入线段树与懒标记机制，将修改操作从单点扩展至区间，在处理 $\mathrm{O}(n)$ 级别的批量更新时避免了逐个节点的冗余操作。题三进一步将标记系统推广至多标记情形——同时维护加法标记与乘法标记——此时标记的合并顺序不再是平凡的，需要精确定义「先乘后加」的线性变换语义以确保正确性。题四将线段树的思想沿时间轴展开，引出可持久化线段树：通过每次修改仅复制 $\mathrm{O}(\log n)$ 个节点，我们得以保留所有历史版本，并借助版本差分技巧在 $\mathrm{O}(\log n)$ 时间内回答静态区间第 $k$ 小这一经典问题。

四道题目从单点到区间、从平凡标记到多标记、从实时到可持久化，每条线索都在前一题的基础上向新的维度延伸，共同勾勒出线段树体系的完整骨架。

报告中的所有代码均使用 C++23 标准编写，在 Apple clang 21.0.0 编译器下通过测试。

\section{题一：树状数组}

\subsection{问题重述}

给定一个长度为 $n$ 的整数序列 $a_1, a_2, \dots, a_n$，需要支持 $q$ 次操作。操作为两种类型：单点加，即将 $a_x$ 的值增加 $k$；区间求和，即查询 $\sum_{i=l}^{r} a_i$ 的值。数据范围 $1 \le n, q \le 10^5$，序列元素及操作数在 32 位有符号整数范围内，下标从 $1$ 开始。

\subsection{$\mathrm{lowbit}$ 与区间划分原理}

处理单点修改与区间查询的最朴素思路是直接维护原数组：修改操作 $\mathrm{O}(1)$ 完成，查询操作 $\mathrm{O}(n)$ 扫描。当 $n$ 与 $q$ 均可达到 $10^5$ 时，最坏情况下 $10^{10}$ 次运算显然不可接受。若反转策略——转而维护前缀和数组——则查询可在 $\mathrm{O}(1)$ 完成，但每次单点修改需要 $\mathrm{O}(n)$ 更新所有受影响的前缀和。两类操作的时间复杂度此消彼长，根源在于朴素数组无法同时兼顾修改与查询的对数时间需求。

树状数组（Fenwick Tree，又称 Binary Indexed Tree）的精妙之处在于，它选取了一组稀疏的前缀和作为节点，使得任意前缀的求和只需合并 $\mathrm{O}(\log n)$ 个节点，且任意单点的修改仅需更新同样数量的节点。这组节点的选取依赖于二进制表示的 $\mathrm{lowbit}$ 性质。

定义 $\mathrm{lowbit}(x) = x \mathbin{\&} (-x)$，即 $x$ 在二进制表示下最低位的 $1$ 所对应的数值。例如，$\mathrm{lowbit}(6) = 2$（因为 $6 = 110_2$，最低位 $1$ 位于第 $2$ 位）。树状数组中的节点 $i$ 维护的是区间 $[i - \mathrm{lowbit}(i) + 1, i]$ 的元素之和。这一设计使得前缀和 $\sum_{j=1}^{x} a_j$ 可以通过反复剥离 $\mathrm{lowbit}$ 来分解：令 $s \gets 0$，不断将树状数组节点 $x$ 的值累加入 $s$，然后令 $x \gets x - \mathrm{lowbit}(x)$，直至 $x = 0$。由于 $\mathrm{lowbit}$ 的值至少为 $1$，每步至少清零一个二进制位，循环次数不超过 $\lfloor\log_2 n\rfloor + 1$，故单次前缀和查询的复杂度为 $\mathrm{O}(\log n)$。

单点修改的机制与之对偶：当 $a_x$ 增加 $k$ 时，所有区间覆盖了位置 $x$ 的树状数组节点都需要同步更新。可以证明，这些节点恰好是 $x, x + \mathrm{lowbit}(x), (x + \mathrm{lowbit}(x)) + \mathrm{lowbit}(x + \mathrm{lowbit}(x)), \dots$，直到超出 $n$ 为止。每次加 $\mathrm{lowbit}$ 至少使当前值增加 $1$，总迭代次数同样不超过 $\mathrm{O}(\log n)$。区间求和 $[l, r]$ 则通过两次前缀和查询的差分实现：$\mathrm{sum}(r) - \mathrm{sum}(l - 1)$。

综上所述，树状数组以 $\mathrm{O}(n)$ 的预处理空间和 $\mathrm{O}(\log n)$ 的单次操作时间，优雅地解决了单点修改与区间查询的矛盾。当 $n = q = 10^5$ 时，约 $2 \times 10^5 \times 17 \approx 3.4 \times 10^6$ 次基本运算在数十毫秒内即可完成。其 $\mathrm{lowbit}$ 原理不仅是实现层面的编码技巧，更深刻地反映了「利用二进制分解将区间操作转化为对数次节点操作」这一数据结构设计的核心思想。

\subsection{代码实现}

\inputminted{c++}{../src/p1_fenwick.cpp}

\section{题二：线段树与区间修改}

\subsection{问题重述}

本题在题一的基础上将修改操作从单点扩展至区间——需要支持对区间 $[l, r]$ 内的每个元素统一增加 $k$，查询操作仍为区间求和。数据范围 $1 \le n, q \le 10^5$，由于区间累加可能导致数值溢出，操作数及中间结果需使用 64 位有符号整数存储，下标从 $1$ 开始。

\subsection{懒标记的延迟传播机制}

题一的树状数组在单点修改上表现优异，但在面对区间修改时却力不从心——若将区间加分解为 $r - l + 1$ 次单点加，则一次区间修改在最坏情况下退化为 $\mathrm{O}(n \log n)$，完全丧失了数据结构的意义。我们需要一种能够将区间修改也压缩至对数时间的数据结构，这便是线段树与懒标记的用武之地。

线段树是一种基于分治思想的二叉树结构。叶节点对应原数组的单个元素，内部节点则对应某个区间内所有元素的聚合信息（本题中为区间和）。对于区间 $[l, r]$，其根节点管辖整个范围，左右子节点分别管辖 $[l, \mathit{mid}]$ 与 $[\mathit{mid}+1, r]$，递归直至叶节点。区间查询同样采用分治策略：若当前节点所管辖的区间完全被查询范围覆盖，则直接返回该节点的聚合值；否则递归进入与查询区间有交集的子节点，并将返回的结果合并。

然而，区间修改若仍沿用此策略——递归进入每一个被覆盖的叶节点并逐一更新——复杂度同样将退化为 $\mathrm{O}(n)$。懒标记正是为解决这一困境而引入的核心机制。其思想极为朴素却深刻：当一次区间修改完整覆盖某个节点的管辖区间时，我们并不立即递归更新其所有后代，而是将该修改信息暂存于一个附加在该节点上的懒标记中，同时仅更新当前节点的聚合值（加上 $k$ 乘以区间长度）。只有当后续的查询或修改操作需要访问该节点的子节点时，才将懒标记下推（pushdown）至两个子节点——即更新子节点的聚合值并为子节点也打上相同的懒标记，随后清除当前节点的懒标记。

这一策略将区间修改的时间复杂度从 $\mathrm{O}(n)$ 降至 $\mathrm{O}(\log n)$：每次修改至多访问 $\mathrm{O}(\log n)$ 个被完整覆盖的节点（这些节点直接接受懒标记），外加 $\mathrm{O}(\log n)$ 个部分覆盖的边界节点（这些节点需要向下递归）。查询操作同样在 $\mathrm{O}(\log n)$ 时间内完成。最终每个操作的平均代价由树高决定，即 $\mathrm{O}(\log n)$。

值得强调的是懒标记的语义：节点 $u$ 上的懒标记 $\mathit{lazy}[u]$ 表示「该节点管辖区间内的每一个元素都尚未被加上的值」。因此，节点聚合值的实时更新遵循 $\mathit{sum}[u] \gets \mathit{sum}[u] + \mathit{lazy}[u] \times \mathit{len}[u]$ 的公式，其中 $\mathit{len}[u]$ 为节点 $u$ 管辖区间的长度。在 pushdown 过程中，父节点的 $\mathit{lazy}$ 被累加到子节点的 $\mathit{lazy}$ 中（支持多次区间加的叠加），子节点的 $\mathit{sum}$ 同样按上述公式更新。线段树节点总数为 $4n$（使用数组存储时的安全上界），建树过程递归遍历所有节点，复杂度 $\mathrm{O}(n)$。$q$ 次操作的总时间复杂度为 $\mathrm{O}((n + q) \log n)$，在 $n = q = 10^5$ 时约 $3.4 \times 10^6$ 次运算。这一延迟传播的设计是线段树高效处理区间修改的基石。

\subsection{代码实现}

\inputminted{c++}{../src/p2_segtree_lazy.cpp}

\section{题三：多标记线段树}

\subsection{问题重述}

在题二的基础上进一步引入区间乘法操作——需要在支持区间加与区间求和的同时，增加对区间 $[l, r]$ 内每个元素乘以 $k$ 的操作。所有区间求和的结果需对模数 $p = 998244353$ 取模后输出。数据范围 $1 \le n, q \le 10^5$，操作数及中间值均在 64 位有符号整数范围内，下标从 $1$ 开始。

\subsection{仿射变换下的标记合并}

题二的线段树仅维护了加法标记，足够应对区间加与区间求和的需求。但当乘法操作加入后，问题的复杂度出现了质的变化——我们同时拥有两种懒标记，而它们之间的合并顺序直接决定了正确性。

设节点 $u$ 当前维护的加法标记为 $\mathit{add}[u]$，乘法标记为 $\mathit{mul}[u]$。它们共同作用于该节点管辖区间内的每个元素 $x$，其语义是：$x$ 的实际值应为 $\mathit{mul}[u] \times x + \mathit{add}[u]$。这是一个典型的仿射变换，其正确性依赖于严格定义的标记合并顺序——乘法必须在加法之前执行。换言之，若先加后乘，则表达式变为 $\mathit{mul}[u] \times (x + \mathit{add}[u]) = \mathit{mul}[u] \times x + \mathit{mul}[u] \times \mathit{add}[u]$，与先乘后加的语义不等价。因此，我们统一将标记语义固定为先乘后加，所有操作都围绕这一约定展开。

当对当前节点施加一个新的乘法标记 $m$ 时：已有标记被更新为 $\mathit{mul}' \gets \mathit{mul} \times m$，$\mathit{add}' \gets \mathit{add} \times m$。原因在于，新的变换 $m \times (\mathit{mul} \times x + \mathit{add}) = (\mathit{mul} \times m) \times x + (\mathit{add} \times m)$，恰好符合先乘后加的形式。当对当前节点施加一个新的加法标记 $a$ 时：$\mathit{add}' \gets \mathit{add} + a$，$\mathit{mul}$ 保持不变，因为 $(\mathit{mul} \times x + \mathit{add}) + a = \mathit{mul} \times x + (\mathit{add} + a)$。

在 pushdown 操作中，父节点的标记 $\mathit{mul}_p$ 和 $\mathit{add}_p$ 被依次施加到子节点上。子节点原有的聚合值 $\mathit{sum}_c$ 被更新为 $\mathit{mul}_p \times \mathit{sum}_c + \mathit{add}_p \times \mathit{len}_c$，子节点的标记同样按上述乘法优先的合并规则更新。随后父节点的标记复位为 $\mathit{mul}_p = 1$、$\mathit{add}_p = 0$，表示没有未下推的修改。

这一标记合并代数的设计将多标记情形统一于仿射变换的框架之下，使得 pushdown 的实现保持了 $\mathrm{O}(1)$ 的常数时间，不影响线段树整体的对数时间复杂度。每次操作仍为 $\mathrm{O}(\log n)$，总复杂度 $\mathrm{O}((n + q) \log n)$，空间复杂度 $\mathrm{O}(n)$。初看之下，多标记的引入似乎大幅增加了实现复杂度，但一旦将每个节点的状态理解为一条直线 $x \mapsto mx + a$，所有操作——修改、pushdown、聚合——都回归为对这条直线的参数的代数操作，逻辑反而变得异常清晰。

\subsection{代码实现}

\inputminted{c++}{../src/p3_segtree_multitag.cpp}

\section{题四：可持久化线段树}

\subsection{问题重述}

给定一个长度为 $n$ 的整数序列 $a_1, a_2, \dots, a_n$，需要回答 $q$ 个询问。每个询问给出三个整数 $l, r, k$，要求回答在子区间 $a_l, a_{l+1}, \dots, a_r$ 中第 $k$ 小的元素。数据范围 $1 \le n, q \le 2 \times 10^5$，$1 \le k \le r - l + 1 \le n$，序列中元素的绝对值不超过 $10^9$，下标从 $1$ 开始。

\subsection{值域离散化与版本差分}

静态区间第 $k$ 小是算法竞赛中的经典问题，其最直观的解法——对每次询问截取子数组并排序——复杂度为 $\mathrm{O}(q \times n \log n)$，在 $n, q = 2 \times 10^5$ 时完全不可行。归并树（Merge Sort Tree）可以在 $\mathrm{O}(\log^2 n)$ 时间内回答单个询问，但常数较大。可持久化线段树（Persistent Segment Tree）提供了时间与空间均为 $\mathrm{O}(n \log n)$ 的优雅解法。算法的核心思路分为两步：值域离散化与版本复用。

首先，注意到原序列中元素的具体数值范围可能很大（本题中可达 $10^9$），但其中真正出现过的不同数值不超过 $n$ 个。因此我们将原序列中的所有数值排序去重，映射到 $[1, m]$ 的连续整数区间上（$m \le n$），这一过程称为离散化。离散化后，问题转化为在 $[1, m]$ 的值域上寻找第 $k$ 小，答案再通过逆映射还原为原值。

其次，对原序列的每一个前缀 $a_1, a_2, \dots, a_i$，我们构建一棵权值线段树，记录该前缀中每种离散化值出现的次数。若对 $n$ 个前缀各构建一棵完整的线段树，空间消耗将高达 $\mathrm{O}(n^2)$，这显然不现实。可持久化的思想在此登场：注意到相邻两棵线段树之间仅有一个元素被插入（$a_i$），从根节点到该元素对应叶节点的路径上共有 $\mathrm{O}(\log n)$ 个节点需要修改，其余节点可以完全复用前一版本的节点。因此，为每个新版本，我们仅需新建 $\mathrm{O}(\log n)$ 个节点，并将未修改的子节点指针直接指向前一版本的对应节点。$n$ 个版本的总节点数为 $\mathrm{O}(n \log n)$，在本题数据规模下约为 $2 \times 10^5 \times 18 \approx 3.6 \times 10^6$，完全在可接受范围内。

回答询问 $(l, r, k)$ 时，我们利用版本差分的思想：版本 $r$ 的权值线段树记录了前 $r$ 个元素的频次分布，版本 $l-1$ 的线段树记录了前 $l-1$ 个元素的频次分布，二者的对应节点相减便得到了区间 $[l, r]$ 内各值的出现次数。查询时，我们从两棵树的根节点同步向下走：检查左子树的频次差值——即区间 $[l, r]$ 中值落在左半值域的元素数量——记为 $\mathit{left\_cnt}$。若 $k \le \mathit{left\_cnt}$，则第 $k$ 小的元素必然在左半值域中，两棵树同时进入左子树；否则该元素在右半值域中，令 $k \gets k - \mathit{left\_cnt}$ 并进入右子树。递归至叶节点时，当前值即为所求。

这一算法的精妙之处在于，它以时间换空间，以可持久化的节点复用策略避免了为每个前缀存储完整线段树，同时利用线段树的二叉树结构将「区间第 $k$ 小」这一看似需要排序才能回答的问题，转化为 $\mathrm{O}(\log n)$ 次比较与计数操作。总体时间复杂度为 $\mathrm{O}((n + q) \log n)$，在 $n = q = 2 \times 10^5$ 时约 $7.2 \times 10^6$ 次运算。空间方面，空线段树占用 $\mathrm{O}(1)$ 个节点，每次插入新增 $\mathrm{O}(\log n)$ 个节点，$n$ 次插入共 $\mathrm{O}(n \log n)$ 个节点，以 $n = 2 \times 10^5$、树高约 $18$ 估算，总内存占用约 $48\,\mathrm{MB}$，在通常的 $256\,\mathrm{MB}$ 限制内绰绰有余。

\subsection{代码实现}

\inputminted{c++}{../src/p4_chairman.cpp}

\section{总结}

本报告以线段树与树状数组为主线，选取了四道代表性题目，系统地展现了区间操作类数据结构从基础到高级的完整图景。

题一的树状数组以 $\mathrm{lowbit}$ 为核心，利用二进制分解将单点修改与区间查询同时压缩至 $\mathrm{O}(\log n)$。其实现之简洁——核心代码不过十余行——与效率之优越形成鲜明对比，使其成为众多算法竞赛选手处理单点修改场景的首选工具。然而，树状数组的局限性也显而易见：它难以优雅地处理区间修改与更复杂的聚合运算，这为线段树的出场埋下了伏笔。

题二引入了线段树与懒标记机制，将修改操作从单点推广至区间，同时保持 $\mathrm{O}(\log n)$ 的时间复杂度。懒标记的核心思想——延迟传播——是数据结构设计中一种普适的优化范式：与其在每次修改时立即更新所有受影响的节点，不如将修改暂存于高层节点，仅在必要时下推。这一策略在各类区间数据结构中均有广泛应用。

题三将懒标记系统从单一加法扩展至加法与乘法并存的情形。多标记引入的核心挑战在于标记之间的合并顺序——不同的顺序导致不同的代数语义，而错误的选择将直接导致数值错误。我们将问题归结为仿射变换 $x \mapsto mx + a$ 的代数，在统一的框架下推导出标记的合并规则。这一「代数化」思维方式提醒我们，在数据结构设计中，将操作形式化为某种代数结构往往能使实现变得清晰且易于验证。

题四将线段树从空间维度扩展至时间维度。可持久化的核心思想——通过节点共享实现版本复用——使得我们得以保存数据结构的所有历史快照而无需承担完整复制的开销。进一步地，版本差分技巧将区间查询转化为两个历史版本的对比，这一策略在可持久化并查集、可持久化平衡树等变体中反复出现，构成了可持久化数据结构的方法论基石。

回顾四道题的递进脉络，可以看到两条贯穿始终的主线。其一，操作范围的扩展——从单点到区间，驱动着数据结构从扁平数组演进为层次化的树形结构。其二，标记系统的丰富——从无标记到单标记、从单标记到多标记代数，每一步都迫使我们重新审视操作的语义并寻找统一的代数表述。这两条线索共同勾勒出线段树体系的演进逻辑：当问题的约束增加时，我们不是在原地修补，而是在更高层次上重新理解数据与操作之间的关系。

\end{document}
